\documentclass[11pt,letterpaper]{article}
% KJO_preamble.tex
% Shared LaTeX preamble for all KJO documentation documents.
% \input this file immediately after \documentclass{}.

% --- Encoding and fonts -------------------------------------------------------
\usepackage[utf8]{inputenc}
\usepackage[T1]{fontenc}
\usepackage{lmodern}
\usepackage{microtype}

% --- Page geometry ------------------------------------------------------------
\usepackage[
  letterpaper,
  top=1in, bottom=1in,
  left=1.25in, right=1in
]{geometry}

% --- Core utilities -----------------------------------------------------------
\usepackage{amsmath}
\usepackage{graphicx}
\usepackage{xcolor}
\usepackage{parskip}        % space between paragraphs; no indent
\usepackage{enumitem}
\usepackage{caption}
\usepackage{subcaption}
\usepackage{float}

% --- Tables -------------------------------------------------------------------
\usepackage{booktabs}
\usepackage{longtable}
\usepackage{array}
\usepackage{multirow}

% --- Code listings ------------------------------------------------------------
\usepackage{listings}
\lstset{
  basicstyle=\ttfamily\small,
  keywordstyle=\bfseries\color{KJOblue},
  commentstyle=\itshape\color{gray},
  stringstyle=\color{KJOgreen!70!black},
  breaklines=true,
  frame=single,
  framesep=4pt,
  xleftmargin=6pt,
  xrightmargin=4pt,
  numbers=left,
  numberstyle=\tiny\color{gray},
  numbersep=8pt,
  language=C++,
}

% --- TikZ / PGF ---------------------------------------------------------------
\usepackage{tikz}
\usetikzlibrary{
  arrows.meta,
  shapes.geometric,
  shapes.misc,
  shapes.symbols,
  positioning,
  fit,
  calc,
  backgrounds,
  decorations.pathreplacing,
  decorations.markings,
  matrix,
}
\usepackage{pgfplots}
\pgfplotsset{compat=1.18}

% --- Headers / footers --------------------------------------------------------
\usepackage{fancyhdr}
\usepackage{lastpage}
\pagestyle{fancy}
\fancyhf{}
\renewcommand{\headrulewidth}{0.4pt}
\renewcommand{\footrulewidth}{0.4pt}
% Per-document: \lhead{}, \rhead{}, \cfoot{} are set in the main .tex file.

% --- Section formatting -------------------------------------------------------
\usepackage{titlesec}
\titleformat{\section}{\large\bfseries\color{KJOblue}}{{\thesection}}{0.8em}{}[\titlerule]
\titleformat{\subsection}{\normalsize\bfseries\color{KJOblue!80!black}}{{\thesubsection}}{0.8em}{}
\titleformat{\subsubsection}{\normalsize\bfseries}{{\thesubsubsection}}{0.8em}{}

% --- Hyperlinks ---------------------------------------------------------------
\usepackage[
  colorlinks=true,
  linkcolor=KJOblue,
  urlcolor=KJOblue!70!black,
  citecolor=KJOblue,
  pdfborder={0 0 0},
]{hyperref}

% --- Color palette -----------------------------------------------------------
\definecolor{KJOblue}   {RGB}{ 30,  90, 160}   % RCU
\definecolor{KJOred}    {RGB}{180,  30,  30}   % EMU
\definecolor{KJOgreen}  {RGB}{ 20, 120,  60}   % GSEMU
\definecolor{KJOorange} {RGB}{200, 110,  20}   % GSE / pressurization
\definecolor{KJOgray}   {RGB}{100, 100, 100}   % Buses / links
\definecolor{KJOpurple} {RGB}{110,  40, 140}   % Shared libraries
\definecolor{KJObg}     {RGB}{248, 248, 252}   % Light background for code

% --- Convenience macros -------------------------------------------------------

% Hardware component name (small caps)
\newcommand{\hw}[1]{\textsc{#1}}

% Software module name (monospace)
\newcommand{\sw}[1]{\texttt{#1}}

% Command name (monospace bold)
\newcommand{\cmd}[1]{\texttt{\textbf{#1}}}

% Pin / signal name
\newcommand{\sig}[1]{\texttt{#1}}

% Note / callout box
\usepackage{tcolorbox}
\tcbuselibrary{skins}
\newtcolorbox{notebox}{
  enhanced,
  colback=KJOblue!5,
  colframe=KJOblue!40,
  fonttitle=\bfseries,
  title=Note,
  left=4pt, right=4pt, top=2pt, bottom=2pt,
}
\newtcolorbox{warningbox}{
  enhanced,
  colback=KJOred!5,
  colframe=KJOred!40,
  fonttitle=\bfseries\color{KJOred},
  title=Warning,
  left=4pt, right=4pt, top=2pt, bottom=2pt,
}

% Title page macro: \KJOtitlepage{title}{subtitle}{version}{date}{description}
\newcommand{\KJOtitlepage}[5]{%
  \begin{titlepage}
    \centering
    \vspace*{2cm}
    {\Huge\bfseries\color{KJOblue} #1 \par}
    \vspace{0.5cm}
    {\Large\color{KJOgray} #2 \par}
    \vspace{1.5cm}
    \rule{0.5\textwidth}{0.4pt}\par
    \vspace{1cm}
    {\large Ken Overton \par}
    \vspace{0.3cm}
    {\normalsize \textit{KJO Liquid Rocket Motor Control System} \par}
    \vspace{1.5cm}
    \begin{tabular}{ll}
      \textbf{Version:} & #3 \\[4pt]
      \textbf{Date:}    & #4 \\
    \end{tabular}
    \vspace{1.5cm}\par
    \begin{minipage}{0.75\textwidth}
      \centering
      \normalsize #5
    \end{minipage}
    \vfill
    {\small\color{KJOgray} \textcopyright\ 2025--2026 Ken Overton.
     All rights reserved. \par}
  \end{titlepage}%
}

\newcommand{\docversion}{0.6}
\newcommand{\docdate}{September 2026}
\lhead{\textsc{KJO Remote Control Unit}}
\rhead{\textsc{RCU Reference} --- v\docversion}
\cfoot{\thepage\ of \pageref{LastPage}}
\hypersetup{pdftitle={KJO RCU Reference}, pdfauthor={Ken Overton}}

\begin{document}
% KJO_tikz_styles.tex
% Shared TikZ node and edge styles for all KJO block diagrams.
% \input this file inside the document body before any tikzpicture environments.

\tikzset{
  % ── Hardware block (peripheral chip / PCB component) ──────────────────────
  hwblock/.style={
    rectangle, draw=KJOgray, line width=1pt,
    fill=KJOgray!10, rounded corners=4pt,
    minimum width=2.6cm, minimum height=0.8cm,
    font=\small, align=center, inner sep=4pt,
  },
  % Sub-variant: MCU (larger, blue border)
  mcublock/.style={
    hwblock,
    draw=KJOblue, fill=KJOblue!8,
    minimum width=3.0cm, minimum height=1.0cm,
    font=\small\bfseries,
  },
  % Sub-variant: power / battery
  powerblock/.style={
    hwblock,
    draw=KJOorange, fill=KJOorange!8,
  },
  % ── Software module box ────────────────────────────────────────────────────
  swmodule/.style={
    rectangle, draw=KJOblue, line width=0.8pt,
    fill=KJOblue!6, rounded corners=3pt,
    minimum width=2.8cm, minimum height=0.7cm,
    font=\small, align=center, inner sep=3pt,
  },
  % Sub-variant: shared library
  sharedlib/.style={
    swmodule,
    draw=KJOpurple, fill=KJOpurple!6,
    dashed,
  },
  % ── Bus / link annotation ──────────────────────────────────────────────────
  busline/.style={
    draw=KJOgray, line width=1.5pt,
    -{Stealth[length=6pt,width=4pt]},
  },
  bidbusline/.style={
    draw=KJOgray, line width=1.5pt,
    {Stealth[length=6pt,width=4pt]}-{Stealth[length=6pt,width=4pt]},
  },
  radioline/.style={
    draw=KJOblue!60, line width=1.5pt, dashed,
    {Stealth[length=6pt,width=4pt]}-{Stealth[length=6pt,width=4pt]},
  },
  gpioline/.style={
    draw=KJOgreen!60!black, line width=1.2pt,
    -{Stealth[length=5pt,width=4pt]},
  },
  % ── Flowchart shapes ──────────────────────────────────────────────────────
  flowproc/.style={
    rectangle, draw=KJOblue!70, line width=0.8pt,
    fill=KJOblue!8, rounded corners=3pt,
    minimum width=3.2cm, minimum height=0.7cm,
    font=\small, align=center, inner sep=3pt,
  },
  flowdec/.style={
    diamond, draw=KJOorange, line width=0.8pt,
    fill=KJOorange!8,
    minimum width=2.2cm, minimum height=0.9cm,
    font=\small, align=center, inner sep=1pt,
  },
  flowterminal/.style={
    rounded rectangle, draw=KJOblue, line width=1pt,
    fill=KJOblue!12,
    minimum width=2.8cm, minimum height=0.7cm,
    font=\small\bfseries, align=center, inner sep=3pt,
  },
  flowarrow/.style={
    draw=KJOgray, line width=0.9pt,
    -{Stealth[length=5pt,width=4pt]},
  },
  % ── Propellant schematic styles ────────────────────────────────────────────
  % Generic pipe line
  pipe/.style={
    draw=black, line width=1.8pt,
  },
  pipearrow/.style={
    pipe,
    -{Stealth[length=7pt,width=5pt]},
  },
  % Oxidizer (N2O) — blue
  oxpipe/.style={
    draw=KJOblue, line width=1.8pt,
  },
  oxpipearrow/.style={
    oxpipe,
    -{Stealth[length=7pt,width=5pt]},
  },
  % Fuel — red
  fuelpipe/.style={
    draw=KJOred, line width=1.8pt,
  },
  fuelpipearrow/.style={
    fuelpipe,
    -{Stealth[length=7pt,width=5pt]},
  },
  % GSE fill line — green
  gsepipe/.style={
    draw=KJOgreen!70!black, line width=1.8pt,
  },
  gsepipearrow/.style={
    gsepipe,
    -{Stealth[length=7pt,width=5pt]},
  },
  % Pressurization — orange dashed
  presspipe/.style={
    draw=KJOorange, line width=1.4pt, dashed,
  },
  presspipearrow/.style={
    presspipe,
    -{Stealth[length=6pt,width=4pt]},
  },
  % Tank (tall rectangle, rounded corners)
  tank/.style={
    rectangle, draw=black, line width=1.5pt,
    rounded corners=4pt, fill=white,
    minimum width=1.8cm, minimum height=2.4cm,
    font=\small\bfseries, align=center, inner sep=4pt,
  },
  n2otank/.style={
    tank, draw=KJOblue, fill=KJOblue!8,
  },
  fueltank/.style={
    tank, draw=KJOred, fill=KJOred!8,
  },
  srctank/.style={
    tank, draw=KJOgreen!70!black, fill=KJOgreen!6,
  },
  % Valve box (servo-actuated)
  valvebox/.style={
    rectangle, draw=black, line width=1pt,
    fill=white, minimum width=2.2cm, minimum height=0.65cm,
    font=\footnotesize\bfseries, align=center, inner sep=3pt,
  },
  gseval/.style={
    valvebox, draw=KJOgreen!70!black, fill=KJOgreen!8,
  },
  emuval/.style={
    valvebox, draw=KJOred, fill=KJOred!6,
  },
  % Component (generic inline component)
  component/.style={
    rectangle, draw=KJOgray, line width=0.8pt,
    fill=white!96!gray, minimum width=2.2cm, minimum height=0.6cm,
    font=\footnotesize, align=center, inner sep=3pt,
  },
  % Instrument callout circle (PT, TT, TC)
  instrument/.style={
    circle, draw=KJOgray, line width=0.8pt,
    fill=white, minimum size=0.6cm,
    font=\tiny\bfseries, align=center, inner sep=0pt,
  },
  % Sensor lead line
  sensorlead/.style={
    draw=KJOgray, line width=0.7pt,
  },
  % Boundary box (dashed outline for GSE / Rocket regions)
  gseboundary/.style={
    draw=KJOgreen!50, line width=1pt, dashed,
    rounded corners=6pt,
  },
  rocketboundary/.style={
    draw=KJOred!40, line width=1pt,
    rounded corners=6pt,
  },
  % Sequence diagram lifeline
  lifeline/.style={
    draw=KJOgray, line width=0.6pt, dashed,
  },
  % Sequence diagram message arrow
  seqarrow/.style={
    draw=KJOblue, line width=0.9pt,
    -{Stealth[length=5pt,width=4pt]},
  },
  % State node (FSM)
  statenode/.style={
    circle, draw=KJOblue, line width=1pt,
    fill=KJOblue!10,
    minimum size=1.2cm,
    font=\small\bfseries, align=center,
  },
  statetransition/.style={
    draw=KJOgray, line width=0.9pt,
    -{Stealth[length=5pt,width=4pt]},
    font=\footnotesize,
  },
}


\KJOtitlepage%
  {Remote Control Unit (RCU)}%
  {Hardware and Software Reference}%
  {\docversion}{\docdate}%
  {This document describes the hardware, software architecture, operator
   interface, and control flow of the Remote Control Unit.  The RCU is the
   sole operator interface to the KJO liquid rocket motor control system.}

\tableofcontents\newpage

%═══════════════════════════════════════════════════════════════════════════════
\section{Hardware Overview}
%═══════════════════════════════════════════════════════════════════════════════

\subsection{Component Summary}

\begin{center}
\begin{tabular}{@{}lll@{}}
\toprule
Component & Part & Interface \\
\midrule
Microcontroller  & Adafruit Feather M0 (ATSAMD21G18A, 48~MHz)   & --- \\
Display          & Adafruit 3.5\textquotedbl\ TFT V2 (HX8357D, 480$\times$320) & SPI \\
Touchscreen      & TSC2007 resistive controller                  & I\textsuperscript{2}C (0x48) \\
LoRa Radio       & HopeRF RFM95W (Semtech SX1276), 915~MHz       & SPI \\
Real-Time Clock  & DS3231 (provides temperature also)            & I\textsuperscript{2}C \\
SD Card          & On Adafruit 3.5\textquotedbl\ TFT shield       & SPI \\
Battery          & 3.7~V LiPo; voltage divider on A7             & ADC A7 \\
\bottomrule
\end{tabular}
\end{center}

\subsection{Pin Assignments}

\begin{center}
\begin{tabularx}{\textwidth}{@{}llY@{}}
\toprule
Signal & Pin & Notes \\
\midrule
\sig{SCREEN\_CS}  & 11  & TFT chip select \\
\sig{SCREEN\_DC}  & 10  & TFT data/command \\
\sig{SD\_CS}      &  5  & SD card chip select \\
\sig{STMPE\_CS}   &  6  & Touch pen-down IRQ (\sig{TOUCH\_IRQ} aliases this pin). Legacy name from an earlier STMPE610-era touch controller; the current TSC2007 is pure I\textsuperscript{2}C and never uses this pin as an SPI chip-select. \\
\sig{RFM95\_CS}   &  8  & Radio chip select \\
\sig{RFM95\_RST}  &  4  & Radio reset \\
\sig{RFM95\_INT}  &  3  & Radio IRQ \\
\sig{BAT\_SENSE}  & A7  & Battery voltage (2:1 divider) \\
\bottomrule
\end{tabularx}
\end{center}

\subsection{Hardware Block Diagram}

\begin{figure}[H]
\centering
\begin{tikzpicture}[node distance=1.1cm and 1.5cm]
  \node[mcublock, minimum width=3.2cm, minimum height=1.2cm] (mcu)
    {Feather M0\\(ATSAMD21G18A)};
  % SPI peripherals (right)
  \node[hwblock, right=2.2cm of mcu, yshift=2.2cm]  (tft)  {HX8357D TFT\\480$\times$320};
  \node[hwblock, right=2.2cm of mcu, yshift=0.8cm]  (sd)   {SD Card};
  \node[hwblock, right=2.2cm of mcu, yshift=-0.6cm] (rf)   {RFM95W LoRa};
  % I2C peripherals (left)
  \node[hwblock, left=2.0cm of mcu, yshift=0.8cm]   (rtc)  {DS3231 RTC};
  \node[hwblock, left=2.0cm of mcu, yshift=-0.4cm]  (touch){TSC2007\\Touch};
  % Power (below)
  \node[powerblock, below=1.3cm of mcu, minimum width=3.0cm] (bat) {LiPo Battery};
  % SPI bus
  \draw[bidbusline] (mcu.east |- tft) -- (tft.west)
    node[midway,above,font=\tiny]{SPI};
  \draw[bidbusline] (mcu.east |- sd)  -- (sd.west)
    node[midway,above,font=\tiny]{SPI};
  \draw[bidbusline] (mcu.east |- rf)  -- (rf.west)
    node[midway,below,font=\tiny]{SPI};
  % I2C bus
  \draw[bidbusline] (mcu.west |- rtc)   -- (rtc.east)
    node[midway,above,font=\tiny]{I\textsuperscript{2}C};
  \draw[bidbusline] (mcu.west |- touch) -- (touch.east)
    node[midway,below,font=\tiny]{I\textsuperscript{2}C};
  % Power
  \draw[busline, draw=KJOorange] (bat.north) -- (mcu.south)
    node[midway,right,font=\tiny]{3.7V};
\end{tikzpicture}
\caption{RCU hardware block diagram.}
\end{figure}

%═══════════════════════════════════════════════════════════════════════════════
\section{Software Architecture}
%═══════════════════════════════════════════════════════════════════════════════

\subsection{Module Summary}

\begin{center}
\begin{tabularx}{\textwidth}{@{}llY@{}}
\toprule
Module & Location & Responsibility \\
\midrule
\sw{main.cpp}          & local  & Setup, loop, command dispatch, GUI logic \\
\sw{KJO\_RCU\_Display} & shared & TFT GUI: buttons, valve symbols, status panels \\
\sw{KJO\_Radio}        & shared & Radio pin constants, frequency, power \\
\sw{KJO\_Command\_Defs}& shared & Command codes, packet struct definitions \\
\sw{KJO\_System\_Config}&shared & Version string, platform identifier \\
\sw{KJO\_Status}       & shared & Tagged Serial/SD logging (\sw{Log\_Message()}); RCU's single-line message panel (\sw{Display\_Message()}) does not use the scrolling \sw{KJO\_Status\_Display} module the other three boards share \\
\bottomrule
\end{tabularx}
\end{center}

\subsection{Main Loop Control Flow}

\begin{figure}[H]
\centering
\begin{tikzpicture}[node distance=0.55cm]
  \node[flowterminal] (start) {loop()};
  \node[flowproc, below=of start]  (rx)  {Check\_and\_Process\\\_Receive\_Packet()};
  \node[flowproc, below=of rx]     (hb)  {Check\_Act\_Heartbeat()};
  \node[flowproc, below=of hb]     (los) {Check\_For\_LOS()};
  \node[flowproc, below=of los]    (vlt) {Check\_Update\_Volts\\\_Link\_Time\_Temp()};
  \node[flowproc, below=of vlt]    (prs) {Check\_Pressure\_Status()};
  \node[flowproc, below=of prs]    (fpg) {Check\_Fill\_Progress\_Status()};
  \node[flowproc, below=of fpg]    (frt) {Check\_Fill\_Cmd\_Retry()};
  \node[flowproc, below=of frt]    (btn) {Check\_Act\_Button\_Press()};
  \node[flowproc, below=of btn]    (clr) {Check\_Clear\_Buttons()};
  \draw[flowarrow] (start)--(rx);
  \draw[flowarrow] (rx)--(hb);
  \draw[flowarrow] (hb)--(los);
  \draw[flowarrow] (los)--(vlt);
  \draw[flowarrow] (vlt)--(prs);
  \draw[flowarrow] (prs)--(fpg);
  \draw[flowarrow] (fpg)--(frt);
  \draw[flowarrow] (frt)--(btn);
  \draw[flowarrow] (btn)--(clr);
  % Loop back
  \draw[flowarrow] (clr.south) -- ++(0,-0.4) -- ++(3.5,0)
    -- ++(0, 11.0) -- (start.east);
\end{tikzpicture}
\caption{RCU main loop control flow.}
\end{figure}

%═══════════════════════════════════════════════════════════════════════════════
\section{Operator Interface (GUI)}
%═══════════════════════════════════════════════════════════════════════════════

\subsection{Button Layout and Behaviour}

Nine buttons are arranged in three rows plus one relocated button, reflecting
the multi-source-ignition redesign (see §\ref{sec:multi-source-ignition}):

\begin{itemize}[nosep]
  \item \textbf{Row 1} (3 equal columns): Fill, Dump, Status.
  \item \textbf{Row 2} (2 columns, original button geometry): Arm/Disarm, Launch.
  \item \textbf{Row 3}: three small ignition-source-select squares (RCU / LCO / AS),
        left-aligned under Arm/Disarm.
  \item \textbf{Abort}: relocated off the button grid entirely -- below the Fuel
        valve indicator, in the lower-right of the screen. Sends the same
        \cmd{RESET} an earlier ``Reset'' button used to send; there is no
        functional distinction between ``Reset'' and ``Abort'' in this system,
        only the relocation and renaming.
\end{itemize}

\begin{sloppypar}
\begin{center}
\begin{tabularx}{\textwidth}{@{}lY@{}}
\toprule
Button & Behaviour \\
\midrule
Fill & Every press sends \cmd{BEGIN\_FILL} and opens the fill-progress popup
       (see §\ref{sec:fill-popup}) -- EMU's own fill-session state decides
       fresh-start vs.\ resume, so RCU doesn't need to track which. Ignored
       once the button shows complete (green); \cmd{RESET} is required
       before filling again. \\
Dump & \textbf{Latching toggle.} First press $\to$ \cmd{DUMP} (button stays
       highlighted). Second press $\to$ \cmd{STOP\_DUMP} (button clears). \\
Status & Sends \cmd{RCU\_STATUS}; momentary highlight, auto-clears. \\
Arm/Disarm & \textbf{Tap-to-toggle.} Label always shows the action a press
       will take. Sends \cmd{ARM\_LAUNCH} (param 1 to arm, 0 to disarm) and
       shows a pending state until EMU acknowledges; the label/highlight
       only commits to Armed/Disarmed on that ack (reverts on failure,
       e.g.\ no ignition source selected). \\
Launch & \textbf{Guarded by confirmation dialog.} Press opens a modal
       CONFIRM / CANCEL overlay; only CONFIRM attempts \cmd{START\_ENGINE},
       and only if still armed for the RCU source at that moment -- see
       §\ref{sec:confirm_dialog}. \\
RCU / LCO / AS & Ignition-source-select squares. Tap sends
       \cmd{SET\_IGNITION\_SOURCE} with the corresponding param (0/1/2);
       shows pending until EMU acknowledges. Rejected by EMU while armed --
       disarm first to change source. \\
Abort & \textbf{Immediate.} Sends \cmd{RESET}; clears all latched button
       states (Fill, Dump) and closes the fill popup if open; brief red
       flash for visual confirmation. \\
\bottomrule
\end{tabularx}
\end{center}
\end{sloppypar}

\begin{notebox}
While armed (\cmd{ARM\_LAUNCH} acknowledged), only Abort, Arm/Disarm, and
Launch (and only if the RCU source is currently armed) respond to touch --
every other button press is silently ignored until disarmed.
\end{notebox}

\subsection{Multi-Source Ignition}
\label{sec:multi-source-ignition}

The engine can be ignited from one of three sources, selected on EMU via
\cmd{SET\_IGNITION\_SOURCE} (param 0=RCU, 1=LCO, 2=Airstart) and rejected by
EMU while armed -- the operator must disarm first to change source. Once a
source is selected, \cmd{ARM\_LAUNCH} arms or disarms it. Only the RCU
source's \cmd{START\_ENGINE} is ever sent from this touchscreen; LCO and
Airstart trigger the ignition sequence through their own independent paths
on EMU/GSEMU, with RCU's role limited to selecting and arming/disarming
the source and observing the unprompted \cmd{ENGINE\_START\_ARMED}
notification EMU sends the instant any source arms the sequence.

\subsection{Fill Progress Popup}
\label{sec:fill-popup}

Pressing Fill opens a modal popup that intercepts all touch except Abort
(which still passes through, so an operator can always abort out of a fill
in progress). While open:

\begin{itemize}[nosep]
  \item RCU polls \cmd{QUERY\_FILL\_PROGRESS} every
        \texttt{FILL\_PROGRESS\_STATUS\_INTERVAL} (1.5~s) -- see §10
        Parameters \& Flags -- updating the Fill button's colour and, if
        the popup is still open, its displayed numbers.
  \item A toggle button reads \textbf{Pause} or \textbf{Resume} depending
        on the last known fill state, sending \cmd{STOP\_FILL} or
        \cmd{BEGIN\_FILL} respectively.
  \item A close control dismisses the popup without changing fill state;
        the fill continues in the background and the Fill button keeps
        reflecting its progress.
  \item If EMU responds \texttt{RELAY\_BUSY} to a Fill/Pause/Resume
        attempt, RCU automatically retries once after a short delay
        (\texttt{FILL\_CMD\_RETRY\_DELAY\_MS}).
\end{itemize}

\subsection{Valve Status Indicators}

Three valve symbols occupy the right portion of the screen:
Fill, N\textsubscript{2}O (Ox), and Fuel.
Each symbol uses a circle-and-bar motif: a vertical bar on a green disc
indicates open; a horizontal bar on a red disc indicates closed.
State is updated from the \cmd{HEARTBEAT} response packet.  There is no
dedicated Dump indicator --- the \cmd{DUMP} command opens the Ox valve
(vents pressurant through the chamber), so its state is reflected on the
Ox indicator.

\begin{sloppypar}
When the tank pressure exceeds \texttt{PRESSURE\_IN\_SYSTEM\_THRESHOLD}
(polled via \cmd{READ\_TANK\_PRESSURE}, \emph{not} carried in the heartbeat
packet), the background of the valve panel turns \textbf{red} as a visual
pressure warning (see §\ref{sec:pressure_warning}).
\end{sloppypar}

\subsection{Status Panels}

\begin{itemize}[nosep]
  \item \textbf{Battery voltages}: RCU voltage (measured locally on A7);
        EMU voltage (received in heartbeat).
  \item \textbf{Link status}: LINK\_OK / LINK\_MARGINAL / LINK\_DOWN, with
        current RSSI value.
  \item \textbf{Time and temperature}: Time from DS3231; temperature
        converted to \textdegree F.
  \item \textbf{QR/CAN status}: a persistent, heartbeat-driven readout
        (\sw{Update\_\allowbreak QR\_\allowbreak Status\_\allowbreak Display()}) showing umbilical connection
        state and GSEMU/LCMU CAN reachability, all three carried in every
        \cmd{HEARTBEAT} response.
  \item \textbf{Message panel}: a single-line status area
        (\sw{Display\_Message()}) showing the most recent status message --
        \emph{not} a scrolling log, and it never displays raw TX/RX
        traffic (that goes to the SD log only, see §9).
\end{itemize}

%═══════════════════════════════════════════════════════════════════════════════
\section{Engine Start Confirmation Dialog}
\label{sec:confirm_dialog}
%═══════════════════════════════════════════════════════════════════════════════

Pressing the \textbf{Launch} button does \emph{not} immediately send
\cmd{START\_ENGINE}.  Instead, a modal overlay is drawn over the main screen
with two large buttons:

\begin{itemize}[nosep]
  \item \textbf{CONFIRM} (green) --- attempts \cmd{START\_ENGINE}, but
        \emph{only} if still armed for the RCU source at the moment of
        touch. If disarmed or re-armed for a different source in the
        meantime (e.g.\ by another operator, or a race with an in-flight
        \cmd{ARM\_LAUNCH}/\cmd{SET\_IGNITION\_SOURCE} ack), the command is
        never sent and the message panel shows ``Launch cancelled -- not
        armed.'' instead.
  \item \textbf{CANCEL} (red) --- dismisses the overlay without transmitting
        any command.
\end{itemize}

Touching anywhere outside both buttons while the overlay is open also
dismisses it (silently, without the ``START cancelled.'' message CANCEL
shows) -- there is no way for the overlay to remain stuck open, and an
accidental touch elsewhere never initiates the ignition sequence.

%═══════════════════════════════════════════════════════════════════════════════
\section{Pressure Warning}
\label{sec:pressure_warning}
%═══════════════════════════════════════════════════════════════════════════════

\begin{sloppypar}
\sw{Check\_Pressure\_Status()} is called every loop iteration, but only
actively polls \cmd{READ\_TANK\_PRESSURE} (every
\texttt{PRESSURE\_STATUS\_INTERVAL}, 2~s) while a fill is active \emph{or}
the warning is already active. Tank pressure is \emph{not} carried in the
heartbeat packet at all -- this poll is RCU's only source for it.
\end{sloppypar}

\begin{warningbox}
Because polling only starts once a fill is active or the warning is
already latched, a pressure rise that occurs outside those two windows
(e.g.\ between fill sessions) is not detected by RCU until the next fill
starts or another mechanism raises it. This is a real monitoring gap in
the current design, not a documentation nuance.
\end{warningbox}

Each \cmd{READ\_TANK\_PRESSURE} response is compared against
\texttt{PRESSURE\_IN\_SYSTEM\_THRESHOLD}. When pressure exceeds it:
\begin{itemize}[nosep]
  \item The valve symbol panel background is redrawn in \textbf{red}
        (suppressed while the fill popup covers that region; repainted
        from the cached flag when the popup closes).
  \item A ``PRESSURISED'' warning label is displayed in the status area.
\end{itemize}

When a subsequent reading returns below the threshold, the panel
background returns to its normal colour and the warning label is cleared.

This visual warning mirrors the EMU-side audible pressure alarm (continuous
buzzer on the EMU) so the operator at the RCU can see that the propellant
system is under pressure, but only while the poll described above is
actually running.

%═══════════════════════════════════════════════════════════════════════════════
\section{Radio Command Dispatch}
%═══════════════════════════════════════════════════════════════════════════════

\sw{RCU\_Execute\_Command(command\_ID, p1)} builds and transmits a command
packet.  \cmd{HEARTBEAT} uses the \texttt{Heartbeat\_Command\_Packet}
structure -- just the 1-byte command code, \emph{not} RCU voltage (a prior
design sent RCU voltage here; it was removed as wasted payload the EMU had
no use for). All other commands use \texttt{Command\_Payload} (4~bytes,
includes a \texttt{short} parameter).

Responses are parsed in \sw{Check\_and\_Process\_Receive\_Packet()}.
The \cmd{HEARTBEAT} response is routed to \sw{Process\_Heartbeat\_Response()},
which updates EMU voltage, the three valve symbols (Fill/Ox/Fuel), and the
QR/GSEMU/LCMU status readout -- all carried in the same packet.

\subsection{Command Reference}

\begingroup
\footnotesize
\begin{longtable}{@{}p{2.6cm}p{4.2cm}p{7.0cm}@{}}
\toprule
Command & Sent by & Notes \\
\midrule
\cmd{HEARTBEAT} & Automatic (every \texttt{heartbeat\_\allowbreak interval}, 5~s, when no other traffic has occurred) & Keeps the link alive; EMU returns voltage, valve positions, and QR/GSEMU/LCMU status \\
\cmd{RESET} & Abort button & Immediate abort; clears all latched button states and closes the fill popup if open \\
\cmd{BEGIN\_\allowbreak FILL} & Fill button (every press); also the fill popup's Resume toggle, and an automatic \texttt{RELAY\_\allowbreak BUSY} retry & Starts/resumes propellant fill; EMU starts pressure recording \\
\cmd{STOP\_\allowbreak FILL} & Fill popup's Pause toggle; automatic \texttt{RELAY\_\allowbreak BUSY} retry & Pauses propellant fill; recording continues until ambient \\
\cmd{QUERY\_\allowbreak FILL\_\allowbreak PROGRESS} & Automatic, every \texttt{FILL\_\allowbreak PROGRESS\_\allowbreak STATUS\_\allowbreak INTERVAL} while a fill is active & Drives the Fill button's colour and the popup's displayed numbers \\
\cmd{READ\_\allowbreak TANK\_\allowbreak PRESSURE} & Automatic, every \texttt{PRESSURE\_\allowbreak STATUS\_\allowbreak INTERVAL}, while a fill is active or the pressure warning is already active & Drives the pressure-warning display (see §\ref{sec:pressure_warning}) -- not a continuous poll \\
\cmd{DUMP} & Dump button (1st press) & Opens Ox valve; vents pressurant through the chamber (no dedicated Dump valve) \\
\cmd{STOP\_\allowbreak DUMP} & Dump button (2nd press) & Closes Ox valve \\
\cmd{SET\_\allowbreak IGNITION\_\allowbreak SOURCE} & RCU/LCO/AS source buttons & Selects the active ignition source (param 0/1/2); rejected by EMU while armed \\
\cmd{ARM\_\allowbreak LAUNCH} & Arm/Disarm button & Arms or disarms the currently-selected source (param 1/0); renamed from \cmd{ENABLE\_\allowbreak LCO\_\allowbreak CONTROL} -- see §11 Version History \\
\cmd{START\_\allowbreak ENGINE} & Launch + CONFIRM (only if still armed for the RCU source) & Arms the 10-step ignition sequence on EMU \\
\cmd{RCU\_\allowbreak STATUS} & Status button & Requests EMU status summary \\
\bottomrule
\end{longtable}
\endgroup

%═══════════════════════════════════════════════════════════════════════════════
\section{SD Logging}
%═══════════════════════════════════════════════════════════════════════════════

Log files are stored in the root of the SD card as \texttt{Log\_N.txt}
(N is auto-incremented to avoid overwriting existing files).
Each entry uses the shared \sw{Log\_Message()} format
(\texttt{"MM/DD/YYYY HH:MM:SS.mmm [TAG] message"}), timestamped from the
DS3231 (falling back to a raw \sw{millis()} count if the RTC never
initialized). TX events use \texttt{TXR}; RX events use \texttt{RXR};
other tags cover boot/status/error events, e.g.:
\begin{lstlisting}[language={},numbers=none]
09/09/2026 10:15:03.142 [TXR] BEGIN_FILL (1) | param: 0
09/09/2026 10:15:03.318 [RXR] BEGIN_FILL (1) | status: OK | r_val: 0 | RSSI: -72
\end{lstlisting}

%═══════════════════════════════════════════════════════════════════════════════
\section{Parameters \& Flags}
%═══════════════════════════════════════════════════════════════════════════════

This section collects the tunable operational parameters and feature flags
compiled into this firmware image, together with the source file where each
is defined. It is meant as a single reference point when tracing or changing
a timing, threshold, or calibration value. Pin/channel wiring assignments,
LoRa command opcodes, and packet byte layouts are documented separately (the
Hardware Overview pin map and the Command Reference sections above) and are
not repeated here. Parameters marked \emph{(shared)} are defined in
\sw{KJO\_\allowbreak Shared\_\allowbreak Libraries} and must agree with the same constant used on
the other board(s) that consume it.

\subsubsection*{Radio Communication}

\begingroup
\footnotesize
\begin{longtable}{@{}p{3.2cm}p{2.4cm}p{6.2cm}p{2.6cm}@{}}
\toprule
Parameter & Value & Description & Source File \\
\midrule
\texttt{RCS\_\allowbreak CENTER\_\allowbreak FREQUENCY} & 917.117~MHz & LoRa carrier frequency; must match on both RCU and EMU radios to communicate. & \texttt{KJO\_\allowbreak Radio.h} (shared) \\
\texttt{TRANSMIT\_\allowbreak POWER} & 5~dBm & RFM95 transmit power (PA\_\allowbreak BOOST pin), valid range 5--23~dBm. & \texttt{KJO\_\allowbreak Radio.h} (shared) \\
\texttt{LINK\_\allowbreak MARGINAL\_\allowbreak THRESHOLD} & $-$90~dBm & RSSI floor below which the link status indicator shows MARGINAL (yellow) instead of OK (green). & \texttt{KJO\_\allowbreak Radio.h} (shared) \\
\texttt{MAX\_\allowbreak PAYLOAD\_\allowbreak LENGTH} & 64~bytes & Size of the LoRa TX/RX payload buffer; sized with headroom above the largest packet (\texttt{Extended\_\allowbreak Heartbeat\_\allowbreak Response\_\allowbreak Packet}, 39 bytes). & \texttt{KJO\_\allowbreak Command\_\allowbreak Defs.h} (shared) \\
\texttt{heartbeat\_\allowbreak interval} & 5000~ms & Interval at which a HEARTBEAT command is sent to the EMU when no other traffic has occurred, to keep the radio link alive. & \texttt{main.cpp} \\
\texttt{LOS\_\allowbreak timeout\_\allowbreak interval} & 30000~ms & Time since last received packet after which the link is declared LOST (LINK\_\allowbreak DOWN) and RSSI/status reset. & \texttt{main.cpp} \\
\texttt{FILL\_\allowbreak CMD\_\allowbreak RETRY\_\allowbreak DELAY\_\allowbreak MS} & 250~ms & Delay before firing the single automatic \cmd{BEGIN\_\allowbreak FILL}/\cmd{STOP\_\allowbreak FILL} retry after EMU responds \texttt{RELAY\_\allowbreak BUSY}. & \texttt{main.cpp} \\
\end{longtable}
\endgroup

\subsubsection*{GUI / Display Timing}

\begingroup
\footnotesize
\begin{longtable}{@{}p{3.2cm}p{2.4cm}p{6.2cm}p{2.6cm}@{}}
\toprule
Parameter & Value & Description & Source File \\
\midrule
\texttt{volt\_\allowbreak RSSI\_\allowbreak update\_\allowbreak interval} & 1000~ms & Refresh interval for the battery voltage, link status, time, and temperature display fields. & \texttt{main.cpp} \\
\texttt{show\_\allowbreak button\_\allowbreak duration} & 1000~ms & How long a momentary button stays highlighted (pressed color) before reverting, for non-latching buttons. & \texttt{main.cpp} \\
\end{longtable}
\endgroup

\subsubsection*{Fill Progress Polling}

\begingroup
\footnotesize
\begin{longtable}{@{}p{3.2cm}p{2.4cm}p{6.2cm}p{2.6cm}@{}}
\toprule
Parameter & Value & Description & Source File \\
\midrule
\texttt{FILL\_\allowbreak PROGRESS\_\allowbreak STATUS\_\allowbreak INTERVAL} & 1500~ms & Poll interval for \cmd{QUERY\_\allowbreak FILL\_\allowbreak PROGRESS} while a fill session is active, driving the Fill button color and popup progress bar. & \texttt{main.cpp} \\
\texttt{PRESSURE\_\allowbreak STATUS\_\allowbreak INTERVAL} & 2000~ms & Poll interval for \cmd{READ\_\allowbreak TANK\_\allowbreak PRESSURE} while a fill is active or the pressure warning is active. & \texttt{main.cpp} \\
\end{longtable}
\endgroup

\subsubsection*{Pressure Monitoring}

\begingroup
\footnotesize
\begin{longtable}{@{}p{3.2cm}p{2.4cm}p{6.2cm}p{2.6cm}@{}}
\toprule
Parameter & Value & Description & Source File \\
\midrule
\texttt{PRESSURE\_\allowbreak IN\_\allowbreak SYSTEM\_\allowbreak THRESHOLD} & 10~PSI & Tank pressure above which the RCU displays the ``system pressurized'' warning (valve panel turns red). & \texttt{KJO\_\allowbreak RCU\_\allowbreak Display.h} (shared) \\
\end{longtable}
\endgroup

\subsubsection*{SD Logging}

\begingroup
\footnotesize
\begin{longtable}{@{}p{3.2cm}p{2.4cm}p{6.2cm}p{2.6cm}@{}}
\toprule
Parameter & Value & Description & Source File \\
\midrule
\texttt{LOG\_\allowbreak FILE\_\allowbreak NAME\_\allowbreak BASE} & \texttt{"Log\_\allowbreak "} & Base filename prefix for SD-card log files; combined with an auto-incrementing index (\texttt{Log\_\allowbreak 0.txt}, \texttt{Log\_\allowbreak 1.txt}, \ldots) to find the next unused file at boot. & \texttt{main.cpp} \\
\end{longtable}
\endgroup

\subsubsection*{Battery Voltage Calibration (RCU Local ADC)}

\begingroup
\footnotesize
\begin{longtable}{@{}p{3.2cm}p{2.4cm}p{6.2cm}p{2.6cm}@{}}
\toprule
Parameter & Value & Description & Source File \\
\midrule
\texttt{RCU\_\allowbreak Batt\_\allowbreak Info.AD\_\allowbreak res} & 12~bits & ADC resolution used when sampling the RCU's own battery voltage divider on pin A7. & \texttt{main.cpp} \\
\texttt{RCU\_\allowbreak Batt\_\allowbreak Info.AD\_\allowbreak samples} & 4 & Number of ADC samples averaged per RCU battery voltage reading (noise reduction). & \texttt{main.cpp} \\
\end{longtable}
\endgroup

%═══════════════════════════════════════════════════════════════════════════════
\section{Version History}
%═══════════════════════════════════════════════════════════════════════════════
\begin{longtable}{@{}p{1.4cm}p{2.2cm}p{11.5cm}@{}}
\toprule
Version & Date & Notes \\
\midrule
0.1 & March 2026 & Initial release. \\
0.2 & March 2026 & Button table updated: FILL and DUMP are latching toggles; Reset is immediate
                   ABORT and clears all latched states; Enable is hold-to-enable LCO (lift = disable);
                   Launch guarded by confirmation dialog.
                   Main loop diagram updated to include \sw{Check\_Pressure\_Status()}.
                   New §4 Engine Start Confirmation Dialog.
                   New §5 Pressure Warning (valve panel background turns red). \\
0.3 & March 2026 & SD log files moved to SD card root (\texttt{Log\_N.txt}); subdirectory path
                   constant (\texttt{LOG\_FILE\_FOLDER}) removed from \sw{main.cpp}. \\
0.4 & July 2026  & Dump valve physically removed from the vehicle. \cmd{DUMP}/\cmd{STOP\_DUMP}
                   retained (still a latching toggle) but now open/close the Ox valve instead of
                   a dedicated Dump valve; the separate Dump valve status indicator was removed
                   and its state folded into the Ox indicator. \\
0.5 & September 2026 & Add new §8 Parameters \& Flags reference section: tunable timing,
                   threshold, and calibration constants across the radio, GUI, fill-progress
                   polling, and pressure-monitoring subsystems, with source-file attribution
                   for each. \\
0.6 & September 2026 & Comprehensive doc-vs-code audit -- this document had fallen well
                   behind two undocumented redesigns already live in the firmware
                   (multi-source ignition, \textasciitilde 2026-07-23; fill-progress popup,
                   \textasciitilde 2026-08-23), neither reflected here until now. Rebuilt \S3
                   Operator Interface around the actual 9-button layout (Fill/Dump/Status,
                   Arm-Disarm/Launch, RCU/LCO/AS source select, a relocated Abort) --
                   replacing a stale 6-button Fill/Dump/Reset/Status/Enable/Launch grid;
                   removed the nonexistent \cmd{ENABLE\_LCO\_CONTROL} command (renamed
                   \cmd{ARM\_LAUNCH} months ago) and documented \cmd{SET\_IGNITION\_SOURCE}/
                   \cmd{ARM\_LAUNCH} and the armed-state lockout (new \S3.2 Multi-Source
                   Ignition); documented the fill-progress popup (new \S3.4) and corrected
                   the Fill button's actual behavior (not a toggle); corrected the Launch
                   confirmation dialog (CONFIRM has a real precondition; touching outside
                   the overlay also silently dismisses it); corrected the Pressure Warning
                   section (tank pressure is never in the heartbeat packet, and polling is
                   \emph{not} continuous -- a real monitoring gap, not just a doc error);
                   rebuilt the Command Reference table (removed the nonexistent command,
                   added 4 missing ones, fixed two mislabeled triggers); corrected the
                   \texttt{Heartbeat\_Command\_Packet} size (1 byte, not 6 -- RCU voltage
                   was removed from it); fixed \S4 Main Loop Control Flow (2 missing
                   steps); fixed the SD Logging example (real timestamp format and tags);
                   corrected the STMPE\_CS pin description (legacy name; it's a touch IRQ,
                   not an SPI chip-select on the current TSC2007 hardware); added
                   \sw{KJO\_Status} to the Module Summary. \\
\bottomrule
\end{longtable}

\end{document}
